\documentclass[11pt]{article}
\usepackage[margin=1in]{geometry}
\usepackage{amsmath,amssymb,amsthm,booktabs,needspace}
\usepackage[hidelinks]{hyperref}
\setlength{\emergencystretch}{2em}
\newtheorem{proposition}{Proposition}
\newtheorem{theorem}[proposition]{Theorem}
\newcommand{\F}{\mathbb F}
\newcommand{\rank}{\operatorname{rank}}
\newcommand{\mr}{\operatorname{mr}}
\title{Betti-gain bounds in every characteristic and exact thirteen-sheet homology}
\author{Internal research working proof}
\date{17 November 2026}
\begin{document}
\maketitle

The thirteen-point test proposed in the preceding ternary
working note is now complete. All $32$ point/dual actions
from its sixteen representation orbits have been retained.
They give sixteen inequivalent covers, two per source knot.
Every cover has rational first Betti number two and exactly
two nontrivial integral torsion factors. The largest Betti
gain is four, below the required thirteen. A separate direct
augmentation calculation handles characteristic thirteen.
No additional target closes. A uniform refinement removes
the prime restriction from the rank implication itself:
the Betti-gain bound holds in every characteristic, although
the split-module equality need not hold.

The all-field exclusion applies only to these induced covers
and their twelve-dimensional augmentation modules. Other
quotients, modules, subgroups and actual cover ranks remain
open. Uniform arguments are written working proofs,
unformalized and externally unreviewed; global novelty is
not claimed.

\section{The group, subgroup and rank gates}

For each source knot let $G=\pi_1(\Sigma_2(K))$. The retained
full $\mathrm{GL}_3(\F_3)$ reflection representations of its
knot group induce epimorphisms
$\rho:G\twoheadrightarrow\mathrm{SL}_3(\F_3)$ of order $5616$.
Determinant is meridian parity, so its kernel is precisely
the image of the double-cover group. Both actions on thirteen
projective points and dual points are transitive. Let
$H=\rho^{-1}(P)$ for a stabilizer in either action.

\begin{proposition}[All-characteristic Betti-gain bound]
For any finitely generated group $G$, subgroup $H$ of finite
index $r\ge2$, and field $k$,
\[
 \rank G\ \ge\ 1+\left\lceil
       \frac{b_1(H;k)-b_1(G;k)}{r-1}\right\rceil.
\]
In particular a gain at least $r$ forces rank at least three,
including characteristics dividing $r$.
\end{proposition}
\begin{proof}
Let $P=k[G/H]$, let $\epsilon:P\to k$ be coordinate sum,
and put $V=\ker\epsilon$, of dimension $r-1$. Transitivity
gives $\dim P^G=1$. Shapiro and the principal-cocycle
dimension therefore give
$\dim Z^1(G,P)=b_1(H;k)+r-1$. The kernel of
$\epsilon_*:Z^1(G,P)\to Z^1(G,k)$ is exactly $Z^1(G,V)$.
Its image has dimension at most $b_1(G;k)$, so
\[
 \dim Z^1(G,V)\ge r-1+b_1(H;k)-b_1(G;k).
\]
For $d$ generators, a cocycle into $V$ is determined by $d$
vectors, so $\dim Z^1(G,V)\le d(r-1)$. Rearranging proves
the assertion. No module splitting or division by $r$ occurs.
\end{proof}

For $r=13$, gain at least thirteen forces cover rank three
over every field. With the retained four-bridge upper and
$\mr(K)\ge1+\rank G$, this proves $\mr(K)=b(K)=4$.
When the characteristic does not divide $r$, the cocycle sum
map is onto and the displayed dimension inequality is an
equality. In a dividing characteristic it may be strict.
The same argument removes the prime restriction from the
earlier index-five rank implication, while retaining the
restriction on its split-module identity. No worldwide
priority claim is made for this elementary consequence.

In characteristic thirteen we also compute $Z^1(G,V)$ directly;
$\dim Z^1\ge25$ would still force $\rank G\ge3$. The invariant
space of $V$ now has dimension one, so principal cocycles
have dimension eleven. The unsplit case is not inferred
from a Betti-number difference.

\section{A positive control exactly at the degree prime}

For every odd prime $p$, take
$G_0=C_p*C_p*C_p=\langle a,b,c\mid a^p,b^p,c^p\rangle$.
Let $a$ act by a $p$-cycle and $b,c$ trivially on $p$ sheets.
The stabilizer preimage has Schreier representatives
$1,a,\ldots,a^{p-1}$. The generator $a^p$ is killed by its
relator; the remaining generators are
$a^iba^{-i},a^ica^{-i}$, with only their order-$p$ relations.
Consequently $H_0=*_{2p}C_p$ and
\[
 b_1(H_0;\F_p)-b_1(G_0;\F_p)=2p-3\ge p.
\]
The all-characteristic bound detects rank three exactly in
the degree characteristic. On the augmentation module all
three norm relations vanish: the full-cycle norm kills
sum-zero vectors, and the other two norms are $pI=0$.
Hence $\dim Z^1(G_0,V)=3(p-1)$, strictly larger than
$(p-1)+(2p-3)=3p-4$ predicted by the split identity.

For $p=13$ an independent integer Smith replay gives
$H_1(H_0;\mathbb Z)=(C_{13})^{26}$. The gain is $23$ and
the direct cocycle dimension is $36$, rather than $35$.
Two altered homology/kernel controls reject. This supports
the uniform proof with a genuine nonsplit positive example;
the all-prime statement is proved by the preceding algebra
and Schreier rewriting, not inferred from this sample.

These are known knot controls. The involutory two-strand
torus presentation gives the dihedral group of order $2p$,
whose meridian-parity kernel is $C_p$. Thus $G_0$ is the
double-cover group of the sum of three $T(2,p)$ knots.
Joining their two-bridge positions gives upper four, and
the new bound gives the matching meridional lower four.
This re-proves known composite cases and claims no new knot.

\section{Exact integral homology}

Write $C_n=\mathbb Z/n\mathbb Z$. Set
\begin{align*}
 \mathcal A&=\{13a_{129},13a_{136},13a_{194},13a_{197}\},\\
 \mathcal B&=\{13a_{135},13a_{137},13a_{196},13a_{198}\}.
\end{align*}
Each of the two induced covers of each knot in one block
has the same homology shown below. Equality of homology
does not identify the manifolds or groups.
\begin{center}
\begin{tabular}{ccc}
\toprule
Block & $H_1(G;\mathbb Z)$ & $H_1(H;\mathbb Z)$\\
\midrule
$\mathcal A$ & $C_{391}$ & $\mathbb Z^2\oplus C_2\oplus C_{882878}$\\
$\mathcal B$ & $C_{401}$ & $\mathbb Z^2\oplus C_2\oplus C_{3128602}$\\
\bottomrule
\end{tabular}
\end{center}
The relevant factorizations are
\[
391=17\cdot23,\quad
882878=2\cdot17\cdot23\cdot1129,\quad
3128602=2\cdot47\cdot83\cdot401.
\]
All displayed odd factors are prime, as checked by bounded
trial division. Thus the gain is four at characteristic two,
three at $1129$ in block $\mathcal A$ and at $47,83$ in block
$\mathcal B$, and two at the other characteristics, including
zero. These gains are stated for $p\ne13$ when used in the
split-module rank gate.

\begin{theorem}
None of the retained point/dual thirteen-sheet covers passes
the augmentation rank-three gate over any field. The maximum
ordinary augmentation cocycle dimension is sixteen, below
twenty-five. In characteristic thirteen it is fourteen.
\end{theorem}
\begin{proof}
Integral Smith homology gives every field Betti number by
the universal coefficient theorem. Away from thirteen,
$\dim Z^1(G,V)=12+b_1(H)-b_1(G)\le16$, with equality at
characteristic two. The separately reconstructed
characteristic-thirteen matrix has $156$ columns, rank $142$
and kernel dimension fourteen. Since $\dim V^G=1$, its first
cohomology dimension is three, while the Betti difference is
two. This exhibits the reason to keep that characteristic
separate. Field extensions preserve all these ranks.
\end{proof}

\section{Representations, kernels and duplicate covers}

The initial data contain two nonconstant reflection
representation orbits per knot. For each orbit retain both
thirteen-point actions. An explicit sheet permutation
intertwines the second orbit's point action with the first
orbit's dual action, and conversely. All generators are
checked, and aliases are permitted only within the same
named source group and normalized presentation.

The first orbit's point and dual actions are inequivalent.
The independent checker enumerates their paired permutation
image of order $5616$. Its point stabilizer has order $432$
and fixes no dual point. An equivariant isomorphism would
have to send the chosen point to such a fixed dual point,
which is impossible. Hence there are two inequivalent
labelled covers per source knot, sixteen in total.

Both projective actions are faithful on the matrix group
$\mathrm{SL}_3(\F_3)$: the independently enumerated permutation
image also has order $5616$. Therefore the intertwiners show
that the two retained representations induce the same
quotient kernel on each source $G$. Among these exact maps
there is one kernel per knot. This neither pools different
knots nor exhausts every possible $\mathrm{SL}_3(\F_3)$
quotient kernel of any $G$.

\section{Construction and independent verification}

The source presentation is the retained double-cover core
presentation \cite{core}, normalized by $x_0=1$:
\[
 \langle x_0,\ldots,x_{12}\mid x_0=1,
   x_ox_a^{-1}x_ox_b^{-1}=1\text{ at each crossing}\rangle.
\]
The producer constructs the lifted edge boundary matrix and
removes twelve edges forming a spanning tree. The resulting
integer matrix has $182$ rows and $157$ columns. Every Smith
reduction retains its elementary unimodular operations.
Base abelianization is computed from the same normalized
presentation.

The independent verifier rebuilds geometric arcs and core
words, constructs projective dual actions by their literal
four-point annihilator planes, and evaluates Fox-prefix
coefficients instead of walking the producer's lifted edges.
It verifies the tree, boundary cancellation and every Smith
operation. The previously shared PD parser supplies only
legacy arc names, checked against geometric paths.

For augmentation coefficients the verifier uses actual
thirteen-coordinate sum-zero vectors. Give each generator
each vector $e_i-e_{12}$ in turn, evaluate the complete affine
word rightmost first at zero, and retain the first twelve
coordinates. This reconstructs every coefficient separately
from the producer's block Fox matrices. Exact modular row
ranks and supplied bases prove kernel completeness at
thirteen; a direct characteristic-two rank also checks the
attained maximum sixteen. All other characteristics are
covered by the Smith factors and the written split-module
argument.

The producer finds maxima for a cyclic base $C_D$ without
an unbounded factorization step: if $p$ first divides Smith
factor $d_j$, it divides every later factor. Its multiplicity
is the number of those factors; subtract one exactly when
$p\mid D$. Repeated gcd removes primes supported on $D$,
and powers of thirteen are removed separately. For these
actual inputs the independent verifier factors the small
integers by trial division and checks every exceptional
prime, plus a generic characteristic.

Two known positive controls use
$G_0=C_3*C_3*C_3$, the triple-trefoil double-cover group.
The first three order-three ternary control matrices give
the point and dual actions. Independent integer homology
replay obtains
\[
 H_1(H_0;\mathbb Z)=\mathbb Z^6\oplus(C_3)^{12},\qquad
 b_1(H_0;\F_3)-b_1(G_0;\F_3)=18-3=15.
\]
Their augmentation cocycle dimensions are $27$ at
characteristic three and $18$ at thirteen, checked directly.
These positive controls detect rank three as intended and
add no new knot case. Four altered Smith-step, free-rank,
characteristic-thirteen kernel and alias-intertwiner inputs
are rejected.

\section{Structural cues and the remaining problem}

The pinned KnotInfo release of 2 November 2026 has the same
Conway polynomial within each block:
\[
 \nabla_{\mathcal A}(z)=-3z^8-7z^6-4z^4+2z^2+1,
 \qquad
 \nabla_{\mathcal B}(z)=3z^8+7z^6+6z^4+4z^2+1.
\]
The retained source subset and archive/CSV hashes document
this observation. The notation fields are blank in that
snapshot, which is not a family exclusion. Neither these
polynomials nor the computed homology prove mutation,
homeomorphic branched covers or common group presentations.

A focused structural question is whether each four-knot block
admits explicit Conway mutations or a common double-cover
presentation. Such a reduction requires geometric moves or
certified group maps before it can reduce the rank problem.
It is not an adopted theorem here.

The completed exclusion concerns these exact covers and
augmentation modules. It leaves other modules, quotient
maps and subgroups available. Actual double-cover ranks,
all original sixteen target equalities, wider priority and
full classical MRC remain open. The retained four-bridge
uppers use the Wirtinger-number theorem \cite{wirt}.

\clearpage
\begin{thebibliography}{9}
\small
\bibitem{core} D. Silver, L. Traldi and S. Williams, Core
groups, J. Knot Theory Ramifications 33 (2024), 2450013.
\url{https://arxiv.org/abs/2311.02048v2}.
\bibitem{wirt} R. Blair, A. Kjuchukova, R. Velazquez and P.
Villanueva, Wirtinger systems of generators of knot groups,
Commun. Anal. Geom. 28 (2020), Theorem 1.3.
\url{https://archive.intlpress.com/site/pub/files/_fulltext/journals/cag/2020/0028/0002/CAG-2020-0028-0002-a002.pdf}.
\bibitem{knotinfo} KnotInfo, pinned database release 2026.11.2.
Exact supplementary fields retained in
\texttt{sources/knotinfo/2026.11.2\_projective13\_geometry\_subset.json}.
Release and CSV hashes are in that subset; this is not a live
classification claim. \url{https://knotinfo.org/}.
\end{thebibliography}
\end{document}
